% ============================================================================
%  Chapter 20 -- Buses and the Video Card                        (Lecture 20)
% ============================================================================
\chapter{Buses and the Video Card}
\label{ch:video}

\epigraphline{\ldots nothing is more serious than gaming\ldots}
  {attributed to Jean Cocteau (1889--1963)}

\begin{chapgoals}
By the end of this chapter you will be able to
\begin{itemize}
\item explain what a peripheral controller is and why devices need synchronisation with the processor;
\item characterise a bus by width, speed and transfer rate, and list its advantages and limits;
\item describe the job of a video card: frame buffer, resolution, colour depth, output interface;
\item compute the memory that a frame needs and the time that it takes to send it over a bus;
\item explain what a GPU does and what general-purpose computing on GPUs is.
\end{itemize}
\end{chapgoals}

% ----------------------------------------------------------------------------
\section{Warm-up on the machine}

\asmfile[caption={One pixel, please: the call, and what the runtime does in principle.},label=lst:vid:pixel]{code/l20/c03.S}

The first four lines of Listing~\ref{lst:vid:pixel} draw one red pixel. The
rest is a sketch of what happens behind the call (it is a sketch of the
principle, not the source of the runtime: \reg{t3} stands for the palette
index that the runtime derives from the colour in \reg{a2}).

\begin{prgconcept}
The screen is an array in memory, with one element per pixel: the
\term{frame buffer}. Drawing a pixel is a store instruction at the address
$\mathit{base}+y\times 320+x$.
\end{prgconcept}

Nothing appears until the frame is sent to the display, which the function
\code{prg32\_gfx\_present} does. On a personal computer these two jobs,
keeping the frame and sending it to the monitor, belong to the video card.
The kit has no video card: its processor does everything, which makes every
step visible. The chapter follows the pixel from the frame buffer to the
display, and then generalises.

% ----------------------------------------------------------------------------
\section{Peripherals and controllers}

\subsection{Input/output devices}

\term{Peripherals} are used for the interaction between the user and the
system. They have limited autonomy: they are controlled and coordinated by
the processor. And they work asynchronously with respect to it, so that
mechanisms of synchronisation are needed: to know when an input datum will be
available, and to know when an output device has finished producing its data.
The terminal, the most common instrument of interaction, is made of two
independent devices, the keyboard and the video; they are not directly
connected, and it is the processor that reproduces on the video what arrives
from the keyboard (Chapter~\ref{ch:io}).

\subsection{Controllers}

\begin{definition}{Controller}
A \term{controller} is the unit that interfaces a peripheral to the bus: it
receives commands from the processor, drives the device, and holds the
registers through which data and state are exchanged.
\end{definition}

Peripherals are interfaced to the processor through dedicated controllers
(Figure~\ref{fig:vid:ctrl}). On the ESP32-C6 the controllers (SPI, I2S, GPIO,
USB) are on the same chip as the processor, and their registers are
memory-mapped (Chapter~\ref{ch:memory}).

\begin{figure}[htb]
\centering
\begin{tikzpicture}
  \node[blkacc,minimum width=24mm,minimum height=9mm] (p) at (0,0) {processor};
  \node[blkvio,minimum width=24mm,minimum height=9mm] (m) at (3.2,0) {memory};
  \draw[prgamber,line width=3pt] (-1.6,-1.5) -- (8.6,-1.5) node[pos=0.86,above,lbl,text=prgamber]{data, address, control};
  \draw[<->,bus] (p) -- (0,-1.5); \draw[<->,bus] (m) -- (3.2,-1.5);
  \foreach \i/\c/\d in {0/{video\\controller}/monitor,1/{USB\\controller}/{keyboard,\\mouse},2/{storage\\controller}/disks}{
    \node[blk,fill=prgsteel,text=white,minimum width=22mm,minimum height=10mm] (c\i) at (\i*3.2+0.6,-3.0) {\c};
    \node[blkgrn,minimum width=22mm,minimum height=9mm] (d\i) at (\i*3.2+0.6,-4.6) {\d};
    \draw[<->,bus] (c\i.north) -- (\i*3.2+0.6,-1.5);
    \draw[<->,thick] (c\i) -- (d\i);
  }
\end{tikzpicture}
\caption{Peripherals reach the bus through their controllers.}
\label{fig:vid:ctrl}
\end{figure}

% ----------------------------------------------------------------------------
\section{Buses}

\subsection{What characterises a bus}

A bus (Chapter~\ref{ch:registers}) is a communication channel shared by
several units, and its characteristics decisively influence the performance
of the whole system. Three numbers describe it.

\begin{description}[style=nextline,leftmargin=1.5em]
\item[Width.] The number of bits of the data bus.
\item[Speed.] The working frequency of the bus; with the width, it gives the
  number of bytes per clock cycle.
\item[Transfer rate.] The megabytes transmitted per second: width times
  frequency.
\end{description}

A bus of 32 bits at 66~MHz, the PCI bus, carries
$4\ \text{bytes}\times 66\cdot 10^{6} = 264$~MB/s. A bus is also a
\emph{standard}: signals, connectors and a protocol on which all the units
agree. The \term{interface} is the hardware module that encodes the data
electrically or optically; the \term{protocol} is the higher-level
description of how the dialogue takes place.

\subsection{Advantages and disadvantages}

With a bus, new devices can easily be added; peripherals can be moved from
one computer to another that uses the same bus; a single set of wires
connects several peripherals; and the cost per byte transmitted is low. On
the other hand the bandwidth of the bus limits the speed of transfer; the
maximum speed depends on the length of the bus and on the number of devices
connected to it; and the bus is a shared resource, on which only one transfer
at a time can take place.

\subsection{Buses through the years}

\begin{table}[htb]
\centering\small
\caption{Buses through the years.}\label{tab:vid:buses}
\begin{tabularx}{\textwidth}{@{}lP{0.23\textwidth}LL@{}}
\toprule
\textbf{Bus} & \textbf{Kind} & \textbf{Used for} & \textbf{Today}\\
\midrule
ISA, EISA, VESA & parallel, 16 to 32 bits & expansion cards of early PCs & disused\\
PCI & parallel, 32 or 64 bits & expansion cards & replaced by PCI Express\\
AGP & parallel, dedicated port & the video card alone & replaced by PCI Express\\
PCI Express & serial, 1 to 16 lanes & video cards, fast storage & current: each generation doubles the speed\\
IDE (ATA), SCSI & parallel & disks & replaced by Serial ATA and NVMe\\
USB & serial & almost every external device & current: from 12 Mbit/s to tens of Gbit/s\\
SPI, I2C, I2S & serial, a few wires & chips on the same board & current: the buses of the \prg kit\\
\bottomrule
\end{tabularx}
\end{table}

The history of Table~\ref{tab:vid:buses} is the passage from parallel to
serial explained in Chapter~\ref{ch:registers}. (The acronyms: ISA,
\emph{industry standard architecture}; PCI, \emph{peripheral component
interconnect}; AGP, \emph{accelerated graphics port}; IDE, \emph{integrated
device electronics}; SCSI, \emph{small computer system interface}; USB,
\emph{universal serial bus}.)

% ----------------------------------------------------------------------------
\section{The video card}

\subsection{What a video card does}

The function of a \term{video card} is to send to the monitor the information
to be displayed. It translates the representation of the image produced by
the processor into a format that the monitor can show. It contains memory of
its own, the \term{video memory}, which holds the frame: the frame buffer.
Every time the screen must be redrawn, it reads the digital representation of
the image and sends the monitor the signal for the colour of each pixel. And
it has processing capability of its own. On \prg the ``video card'' is a
region of RAM plus the SPI peripheral: no separate hardware at all.

\subsection{Analogue and digital monitors}

For a CRT monitor, which has an analogue input, the card converts the digital
image into analogue signals; the conversion is done by a component called
RAMDAC (\emph{random access memory digital-analog converter}). The VGA output
of 1987 was designed for cathode-ray tubes, and suffers from electrical noise,
distortion and errors in the placing of the pixels. For an LCD monitor, which
has a digital input, the card adapts the colours and the resolution of the
image, provides local memory to use as a buffer and supports graphics
acceleration, for example in three dimensions; with DVI each pixel of the
output corresponds to one pixel of the screen, at its native resolution.

\begin{table}[htb]
\centering\small
\caption{Video outputs.}\label{tab:vid:outputs}
\begin{tabularx}{\textwidth}{@{}lllL@{}}
\toprule
\textbf{Interface} & \textbf{Year} & \textbf{Signal} & \textbf{Notes}\\
\midrule
VGA & 1987 & analogue & for CRT monitors; kept for years for compatibility\\
DVI & 1999 & digital (and analogue) & introduced with LCD monitors\\
HDMI & 2003 & digital & replaces the earlier standards; carries audio too\\
DisplayPort & 2006 & digital & video, audio and other data; also over USB-C connectors\\
SPI to a display controller & -- & digital, serial & the \prg kit: an ILI9341 panel of 320 by 240 pixels\\
\bottomrule
\end{tabularx}
\end{table}

The ILI9341 of the last row of Table~\ref{tab:vid:outputs} is the chip on the
display module of the kit: it has its own memory for the frame and drives the
liquid crystals (Chapter~\ref{ch:io}).

\subsection{Resolution: the size of the pixel grid}

The \term{resolution} determines the definition of the image, and how many
pixels must be stored and sent for every frame (Table~\ref{tab:vid:res}).
The traditional ratio of 4:3 has given way to the panoramic 16:9.

\begin{table}[htb]
\centering\small
\caption{Some resolutions.}\label{tab:vid:res}
\begin{tabular}{@{}lrll@{}}
\toprule
\textbf{Resolution} & \textbf{Pixels} & \textbf{Aspect ratio} & \textbf{Where}\\
\midrule
$320\times 200$ & 64\,000 & 8:5 & Commodore 64, VGA games, the \prg viewport\\
$640\times 480$ & 307\,200 & 4:3 & VGA\\
$1024\times 768$ & 786\,432 & 4:3 & monitors of the 1990s\\
$1920\times 1080$ & 2\,073\,600 & 16:9 & Full HD\\
$3840\times 2160$ & 8\,294\,400 & 16:9 & 4K\\
\bottomrule
\end{tabular}
\end{table}

\subsection{Colour depth: bits per pixel}

The \term{colour depth} is the number of bits used for each pixel
(Table~\ref{tab:vid:depth}). In True Color every colour is obtained from 256
shades of each of the three additive primaries, red, green and blue. In the
mode with 8 bits the value of a pixel is not a colour but an \emph{index}
into a table of 256 colours, the \term{palette}.

\begin{table}[htb]
\centering\small
\caption{Colour depths.}\label{tab:vid:depth}
\begin{tabular}{@{}lrll@{}}
\toprule
\textbf{Depth} & \textbf{Colours} & \textbf{Bytes per pixel} & \textbf{Name}\\
\midrule
4 bits & 16 & 0.5 & standard VGA\\
8 bits & 256 & 1 & 256-colour mode: an index into a palette\\
16 bits & 65\,536 & 2 & High Color: RGB565\\
24 bits & 16\,777\,216 & 3 & True Color\\
\bottomrule
\end{tabular}
\end{table}

\prg uses both middle rows: a frame buffer of 8 bits per pixel that holds
palette indices, and a palette of 256 colours in RGB565
(Chapter~\ref{ch:information}).

\begin{theorem}{Memory of a frame}
A frame of $W\times H$ pixels with a depth of $d$ bits per pixel occupies
$W\cdot H\cdot d/8$ bytes.
\end{theorem}

For instance $640\times 480$ in True Color needs
$307\,200\times 3 = 921\,600$ bytes.

\subsection{Following the pixel on PRG32}

\begin{figure}[htb]
\centering
\begin{tikzpicture}[node distance=4mm]
  \node[blkacc,text width=21mm,minimum height=16mm] (a) {\textbf{your code}\\[1pt]\scriptsize\texttt{prg32\_gfx\_pixel}};
  \node[blkvio,text width=21mm,minimum height=16mm,right=of a] (b) {\textbf{frame buffer}\\[1pt]\footnotesize $320\times 200$ bytes in RAM};
  \node[blkamb,text width=21mm,minimum height=16mm,right=of b] (c) {\textbf{palette}\\[1pt]\footnotesize 256 entries of RGB565};
  \node[blkred,text width=21mm,minimum height=16mm,right=of c] (d) {\textbf{SPI bus}\\[1pt]\footnotesize 27 MHz, serial, with DMA};
  \node[blkgrn,text width=21mm,minimum height=16mm,right=of d] (e) {\textbf{display}\\[1pt]\footnotesize ILI9341, $320\times 240$};
  \foreach \x/\y in {a/b,b/c,c/d,d/e} \draw[->,thick] (\x) -- (\y);
\end{tikzpicture}
\caption{The journey of a pixel on the \prg kit.}
\label{fig:vid:journey}
\end{figure}

Figure~\ref{fig:vid:journey} is the journey of the pixel of the warm-up, and
Table~\ref{tab:vid:numbers} gives its numbers (from
\fname{docs/hardware/ili9341.md} of the repository).

\begin{table}[htb]
\centering\small
\caption{The arithmetic of a frame on \prg.}\label{tab:vid:numbers}
\begin{tabularx}{\textwidth}{@{}LL@{}}
\toprule
\textbf{Quantity} & \textbf{Value}\\
\midrule
memory for one frame, one byte per pixel & $320\times 200\times 1 = 64\,000$ bytes\\
the same frame at two bytes per pixel & 128\,000 bytes: a quarter of all the RAM of the chip\\
bits sent to the display for a whole frame & $64\,000$ pixels $\times$ 16 bits $=1\,024\,000$ bits\\
time on the bus at 27 MHz & about 38 ms: at most some 26 full frames per second\\
\bottomrule
\end{tabularx}
\end{table}

The last row follows from $1\,024\,000/27\,000\,000 = 37.9$~ms. The platform
gives two engineering answers to these numbers: an \emph{indexed} frame
buffer, which halves the memory, and the tracking of the rows that have
changed, so that only those are sent. Chapter~\ref{ch:verm1p4} will exploit
the second.

\subsection{Memory and bandwidth: the pressure on video hardware}

The memory and the processor time needed by ever more demanding graphics,
three-dimensional video games above all, imposed a rapid evolution on video
cards. As the definition of the images grows, so does the memory needed to
represent a frame, and the flow of data from the memory to the video card
risks saturating the capacity of the bus: a 4K frame in True Color at 60
frames per second is $8.3\cdot 10^{6}\times 3\times 60$ bytes, about 1.5~GB
of pixels every second. There were two lines of attack: more and more memory
on the card itself, and a wider or dedicated bus, AGP yesterday and PCI
Express today. Several cards could even be connected together to share the
work; today a single card, or many cards in a server, are the rule.

% ----------------------------------------------------------------------------
\section{GPU and GPGPU}

\subsection{The graphics processing unit}

The share of processor time spent computing images grew with the use of
graphical interfaces. To relieve the CPU, video cards were given dedicated
processors.

\begin{definition}{GPU}
A \term{graphics processing unit} (GPU) is a processor specialised in the
operations of graphics, which it executes in parallel on many pixels.
\end{definition}

A GPU accelerates raster graphics, that is, the copying of blocks of pixels
(\emph{blitting}) and the operations on textures, and three-dimensional
graphics: transformations, texture mapping, lighting. Programs use it through
libraries: OpenGL, Direct3D, and today Vulkan and Metal. A \term{device
driver} adapts the software that runs on the CPU to the primitives offered by
the GPU; what is not supported is emulated by the CPU, slowly.

\begin{prgmachine}
On \prg there is no GPU: \code{prg32\_gfx\_rect} is a loop of store
instructions executed by the \riscv core. This is why the frame rate of a
busy cartridge drops, and why the cost of every pixel can be seen.
\end{prgmachine}

\subsection{General-purpose computing on GPUs}

Modern GPUs execute the same operation in parallel on whole matrices of
pixels, with a very large number of simple cores: hundreds, thousands.
\term{GPGPU} is the use of GPUs for applications of a general kind, not
necessarily graphical (Figure~\ref{fig:info:gpgpu} showed the simplest
example, a parallel sum). The programming models are CUDA (\emph{Compute
Unified Device Architecture}, by NVIDIA) and OpenCL, an open standard.
Scientific simulation, data analysis and, above all today, the training and
the use of neural networks run on GPUs; so do the weather and marine models
of the Department. Chapter~\ref{ch:beyond} returns to the subject.

% ----------------------------------------------------------------------------
% Generated by tools/textbook/export_listings.js from decks/plan.js and code/l20.S: do not edit.
\section{The program of this chapter}\label{sec:lab:l20}

The listings of this chapter are parts of one program, \fname{code/l20.S} in the
repository of the course, printed in full in Listing~\ref{lst:lab:l20}. The
automatic tests of the course material build it and run it on the headless
runner of \prgqt{}, with scripted input, and check the frame;
Figure~\ref{fig:lab:l20} is the frame that the test wrote.

\begin{description}[style=nextline,leftmargin=1.5em]
\item[On the screen.] A single red pixel at x = 100, y = 50 on a black screen. Look closely.
\item[What to press.] Nothing: the program runs by itself.
\end{description}

\begin{figure}[htb]
\centering
\includegraphics[width=0.5\textwidth]{lab_l20.png}
\caption{\fname{l20.S} on \prgqt{}.}\label{fig:lab:l20}
\end{figure}

\asmfile[caption={The complete program of the chapter, \fname{code/l20.S}.},label=lst:lab:l20]{code/l20/program.S}

To run it on \prgqt{} the commands are those of Listing~\ref{lst:lab:l20:run};
Section~\ref{sec:prg:testing} gives the same steps for QEMU and for the board.

\begin{shell}[caption={Building \fname{l20.S} and running it on \prgqt{}.},label=lst:lab:l20:run]
# 1. start PRG32-QT on this computer (it listens on port 8080)

# 2. in the PRG32 folder, with ESP-IDF active: build the cartridge
python3 -m prg32 cartridge build ~/arc12/code/l20.S --portable \
    --entry-prefix l20 --name l20 --out l20.prg32

# 3. send it to PRG32-QT and start it
python3 -m prg32 esp32c6 upload-and-run l20.prg32 \
    --url http://localhost:8080 --slot cart0

# 4. look, press, compare with the previous slide
\end{shell}


% ----------------------------------------------------------------------------
\section{Check your understanding}

\begin{quiz}{How much memory does a frame of $640\times 480$ pixels need in True Color (24 bits)?}
\item 307\,200 bytes
\item 614\,400 bytes
\item 921\,600 bytes
\item 7\,372\,800 bytes
\end{quiz}

\begin{quiz}{A 16-bit bus works at 50 MHz, one transfer per cycle. What is its transfer rate?}
\item 50 MB/s
\item 100 MB/s
\item 400 MB/s
\item 800 MB/s
\end{quiz}

\begin{quiz}{In an 8-bit indexed frame buffer, the byte of a pixel contains:}
\item its red component
\item its brightness
\item the position of its colour in a palette
\item its $x$ coordinate
\end{quiz}

\begin{quiz}{What is GPGPU?}
\item A faster video output
\item Using graphics processors for non-graphical computation
\item A bus for video cards
\item A kind of video memory
\end{quiz}

\section{Exercises}

\begin{exercise}[paper first, then \prg]
\begin{enumerate}[nosep]
\item Compute the memory for one 4K frame in True Color, and the bytes per
  second at 60 frames per second.
\item The display bus of \prg runs at 27~MHz: how long does it take to send 20
  rows of 320 pixels at 16 bits per pixel?
\item Draw a diagonal line pixel by pixel, with \code{prg32\_gfx\_pixel} in a
  loop.
\item Fill the screen once with \code{prg32\_gfx\_clear} and once with 64\,000
  calls to \code{prg32\_gfx\_pixel}: compare the frame rates.
\item Change one entry of the palette with \code{prg32\_palette\_set} at every
  frame, and watch every pixel that uses it change at once.
\end{enumerate}
\expected{On paper: 24\,883\,200 bytes per frame, about 1.5~GB per second;
102\,400 bits, about 3.8~ms. On \prg: palette cycling animates the screen
without redrawing a single pixel. It is the reason why indexed colour was
loved by the programmers of games: one store changes thousands of pixels.}
\end{exercise}

\section{Glossary}

\begin{glossarytable}
peripheral & periferica\\
controller & controller (unit\`a di controllo della periferica)\\
bus width & larghezza del bus\\
transfer rate & velocit\`a di trasferimento\\
video card & scheda video\\
resolution & risoluzione\\
colour depth & profondit\`a di colore\\
palette & tavolozza (palette)\\
video memory & memoria video\\
device driver & driver di periferica\\
\end{glossarytable}

\section{Conclusions and reading}

Input/output devices are a fundamental component of computers, and each is
managed by a controller that interfaces it to the bus. A bus is characterised
by width, speed and transfer rate, and its bandwidth is a limit to design
round. Video cards make it possible to display graphics and text: a frame
buffer, a resolution, a colour depth, an output. Their processors, the GPUs,
are also used for high-performance computing.

\paragraph{Reading.} \citet{tanenbaum2013}, chapters 2 and~3 (buses,
displays); \emph{Playing with RISC-V}, chapter~7, the graphics game
\citep{prg32textbook}; \fname{docs/hardware/ili9341.md} in the \prg
repository.
